% ============================================================
%  Chapter 3: Oscillatory Necessity and the Discrete Spectrum
%  Source: bounded-phase-space-trajectory.tex, Consequences 1--5
% ============================================================

\epigraph{%
  Boundedness forces return.\\
  Return forces oscillation.\\
  Oscillation forces discreteness.%
}{}

\section{Poincaré Recurrence}
\label{sec:ch3-poincare}

The first consequence of Axiom~\ref{ax:bpsl} is the oldest result in
ergodic theory.

\begin{consequence}{Poincaré Recurrence}{poincare}
  Let $(\Omega, \mathcal{B}, \mu, \phi_t)$ satisfy BPSL\@.
  Then for any measurable set $A \subset \Omega$ with $\mu(A) > 0$,
  $\mu$-almost every $x \in A$ satisfies $\phi_{t_n}(x) \in A$ for an
  infinite sequence $t_1 < t_2 < \cdots \to \infty$.
\end{consequence}

\begin{proof}[Proof sketch]
  Standard: apply the Poincaré recurrence theorem to the measure-preserving
  flow $\phi_t$ on the finite-measure space $(\Omega, \mu)$.
  See~\cite{Walters1982} for the complete argument.
\end{proof}

\begin{remark}
  Recurrence is the formal statement that the system does not escape.
  Combined with the fact that any trajectory must move (it cannot be static
  in a generic bounded domain), recurrence forces the trajectory to loop.
\end{remark}

\section{Non-Persistence Without Boundedness}
\label{sec:ch3-nopersist}

\begin{consequence}{Non-Persistence of Unbounded Systems}{noboundsnopersistence}
  If $\mu(\Omega) = \infty$, a measure-preserving flow need not be recurrent,
  and systems generically drift to infinity, contradicting persistence.
\end{consequence}

This consequence has immediate physical content: any system observed to be
persistent is bounded.
The mass spectrometer exists precisely because ions in transit are transiently
bounded by the instrument's fields.

\section{Oscillatory Necessity}
\label{sec:ch3-oscillatory}

Let $x(t)$ be a trajectory of the system.
We say the trajectory is \emph{oscillatory} if $x(t)$ is not eventually
monotone and not eventually static.

\begin{consequence}{Oscillatory Necessity}{oscillatory}
  Every trajectory of a persistent system (one satisfying BPSL) is oscillatory.
\end{consequence}

\begin{proof}[Proof sketch]
  Three Lemmas eliminate the alternatives:
  \begin{itemize}
    \item \textit{Lemma (Static elimination)}: a static trajectory has
      zero-measure initial conditions and is excluded by the measure-preserving
      hypothesis and generic arguments.
    \item \textit{Lemma (Monotone elimination)}: a trajectory that is
      eventually monotone escapes to the boundary of $\Omega$, contradicting
      $\mu(\Omega) < \infty$.
    \item \textit{Lemma (Chaotic non-recurrence)}: a chaotic trajectory that
      does not return to any neighbourhood of its initial condition contradicts
      Consequence~\ref{cons:poincare}.
  \end{itemize}
  Oscillatory motion is the only remaining alternative.
\end{proof}

\section{The Discrete Koopman Spectrum}
\label{sec:ch3-koopman}

Define the Koopman operator $U_t : \Hilbert \to \Hilbert$ by
$U_t f(x) = f(\phi_t(x))$, where $\Hilbert = L^2(\Omega, \mu)$.
Clause~(ii) of BPSL makes $U_t$ unitary.
Define the generator $H$ by $U_t = e^{-itH}$.

\begin{consequence}{Discrete Koopman Spectrum}{discretemodes}
  If $\Omega$ is compact (a consequence of clause~(i) when $\Omega$ is a
  manifold), then $H$ has purely discrete spectrum:
  $\sigma(H) = \{\omega_k\}_{k\in\mathbb{N}}$
  with each $\omega_k$ an eigenvalue of finite multiplicity.
\end{consequence}

\begin{proof}[Proof sketch]
  Compactness of $\Omega$ implies that $U_t$ is almost periodic (Bochner's
  theorem).
  An almost periodic unitary group has purely discrete spectrum.
  Alternatively: $\phi_t$ on compact $\Omega$ admits an ergodic decomposition
  into tori, and the Fourier modes on each torus form the discrete spectrum.
\end{proof}

The discrete spectrum $\{\omega_k\}$ is the frequency basis of the oscillatory
dynamics.
In quantum mechanics, these are the transition frequencies; in classical
mechanics, the Fourier components of the periodic trajectory.

\section{Energy Quantisation: $E = \hbar\omega$}
\label{sec:ch3-equantum}

\begin{consequence}{Energy Quantisation}{energyfrequency}
  For each Koopman eigenfrequency $\omega_k$, the action of the corresponding
  mode is quantised in units of $\hbar$:
  \begin{equation}\label{eq:ehbaromega}
    E_k = \hbar\omega_k.
  \end{equation}
\end{consequence}

\begin{proof}[Proof sketch]
  The Bohr--Sommerfeld quantisation condition $\oint p\,dq = nh$ applied to a
  periodic orbit in a bounded phase space gives $E = nh\nu = \hbar\omega_n$.
  The Koopman eigenfrequencies are the frequencies of these periodic orbits.
  Thus $E_k = \hbar\omega_k$ without postulating the Planck relation.
\end{proof}

\begin{remark}
  Planck's relation $E = h\nu$ is usually taken as a fundamental postulate.
  Here it is a theorem: the consequence of applying Bohr--Sommerfeld
  quantisation to the discrete Koopman spectrum of a bounded system.
  The Planck constant $\hbar$ enters as the quantum of action, itself
  derivable from the minimal phase-space cell area.
\end{remark}

\bigskip
\begin{tcolorbox}[colback=crownblue!5, colframe=crownblue!60!black,
                  title={\textbf{Chapter 3 Summary}}]
\begin{itemize}[noitemsep]
  \item Boundedness implies Poincaré recurrence (Consequence~\ref{cons:poincare}).
  \item Recurrence + non-stasis + non-monotonicity $\Rightarrow$ oscillatory
    dynamics (Consequence~\ref{cons:oscillatory}).
  \item Oscillations in a compact bounded domain have a discrete frequency
    spectrum (Consequence~\ref{cons:discretemodes}).
  \item Each mode satisfies $E_k = \hbar\omega_k$ by Bohr--Sommerfeld
    quantisation (Consequence~\ref{cons:energyfrequency}).
  \item No quantum postulates were used.  Quantisation follows from boundedness.
\end{itemize}
\end{tcolorbox}
